\documentclass[10pt,a4paper]{amsart}
\usepackage{amsmath,amssymb,amsthm,graphicx}
\usepackage{verbatim}
\newenvironment{reference}{\comment}{\endcomment}
%\newenvironment{reference}{}{}
\newenvironment{slogan}{\comment}{\endcomment}
\newenvironment{history}{\comment}{\endcomment}

% For commutative diagrams we use Xy-pic
\usepackage[all]{xy}

% We use 2cell for 2-commutative diagrams.
\xyoption{2cell}
% Theorem environments.
%
\theoremstyle{plain}
\newtheorem{theorem}[subsection]{Theorem}
\newtheorem{proposition}[subsection]{Proposition}
\newtheorem{lemma}[subsection]{Lemma}

\theoremstyle{definition}
\newtheorem{definition}[subsection]{Definition}
\newtheorem{example}[subsection]{Example}
\newtheorem{exercise}[subsection]{Exercise}
\newtheorem{situation}[subsection]{Situation}

\theoremstyle{remark}
\newtheorem{remark}[subsection]{Remark}
\newtheorem{remarks}[subsection]{Remarks}

\numberwithin{equation}{subsection}

% Macros
%
\def\lim{\mathop{\mathrm{lim}}\nolimits}
\def\colim{\mathop{\mathrm{colim}}\nolimits}
\def\Spec{\mathop{\mathrm{Spec}}}
\def\Hom{\mathop{\mathrm{Hom}}\nolimits}
\def\Ext{\mathop{\mathrm{Ext}}\nolimits}
\def\SheafHom{\mathop{\mathcal{H}\!\mathit{om}}\nolimits}
\def\SheafExt{\mathop{\mathcal{E}\!\mathit{xt}}\nolimits}
\def\Sch{\mathit{Sch}}
\def\Mor{\mathop{\mathrm{Mor}}\nolimits}
\def\Ob{\mathop{\mathrm{Ob}}\nolimits}
\def\Sh{\mathop{\mathit{Sh}}\nolimits}
\def\NL{\mathop{N\!L}\nolimits}
\def\CH{\mathop{\mathrm{CH}}\nolimits}
\def\proetale{{pro\text{-}\acute{e}tale}}
\def\etale{{\acute{e}tale}}
\def\QCoh{\mathit{QCoh}}
\def\Ker{\mathop{\mathrm{Ker}}}
\def\Im{\mathop{\mathrm{Im}}}
\def\Coker{\mathop{\mathrm{Coker}}}
\def\Coim{\mathop{\mathrm{Coim}}}

% Boxtimes
%
\DeclareMathSymbol{\boxtimes}{\mathbin}{AMSa}{"02}

%
% Macros for moduli stacks/spaces
%
\def\QCohstack{\mathcal{QC}\!\mathit{oh}}
\def\Cohstack{\mathcal{C}\!\mathit{oh}}
\def\Spacesstack{\mathcal{S}\!\mathit{paces}}
\def\Quotfunctor{\mathrm{Quot}}
\def\Hilbfunctor{\mathrm{Hilb}}
\def\Curvesstack{\mathcal{C}\!\mathit{urves}}
\def\Polarizedstack{\mathcal{P}\!\mathit{olarized}}
\def\Complexesstack{\mathcal{C}\!\mathit{omplexes}}
% \Pic is the operator that assigns to X its picard group, usage \Pic(X)
% \Picardstack_{X/B} denotes the Picard stack of X over B
% \Picardfunctor_{X/B} denotes the Picard functor of X over B
\def\Pic{\mathop{\mathrm{Pic}}\nolimits}
\def\Picardstack{\mathcal{P}\!\mathit{ic}}
\def\Picardfunctor{\mathrm{Pic}}
\def\Deformationcategory{\mathcal{D}\!\mathit{ef}}
\usepackage{cleveref}
\crefname{originallemma}{original lemma}{original lemmas}
\Crefname{originallemma}{Original lemma}{Original lemmas}
\makeatletter
\newlabel{derived-lemma-fully-faithful-adjoint-kernel-zero}{{09J1}{source}{}{sourcepointer-1}{}}
\newlabel{derived-lemma-fully-faithful-adjoint-kernel-zero@cref}{{[originallemma][0][]09J1}{[1][0][]source}{}{}{}}
\makeatother
\crefname{originallemma}{original lemma}{original lemmas}
\Crefname{originallemma}{Original lemma}{Original lemmas}
\makeatletter
\newlabel{dga-lemma-fully-faithful-in-compact-case}{{09LZ}{source}{}{sourcepointer-2}{}}
\newlabel{dga-lemma-fully-faithful-in-compact-case@cref}{{[originallemma][0][]09LZ}{[1][0][]source}{}{}{}}
\makeatother
\crefname{originallemma}{original lemma}{original lemmas}
\Crefname{originallemma}{Original lemma}{Original lemmas}
\makeatletter
\newlabel{dga-lemma-tensor-with-complex-derived}{{09LX}{source}{}{sourcepointer-3}{}}
\newlabel{dga-lemma-tensor-with-complex-derived@cref}{{[originallemma][0][]09LX}{[1][0][]source}{}{}{}}
\makeatother
\crefname{originallemma}{original lemma}{original lemmas}
\Crefname{originallemma}{Original lemma}{Original lemmas}
\makeatletter
\newlabel{dga-lemma-tensor-with-complex-hom-adjoint}{{09LY}{source}{}{sourcepointer-4}{}}
\newlabel{dga-lemma-tensor-with-complex-hom-adjoint@cref}{{[originallemma][0][]09LY}{[1][0][]source}{}{}{}}
\makeatother
\crefname{originalsection}{original section}{original sections}
\Crefname{originalsection}{Original section}{Original sections}
\makeatletter
\newlabel{dga-section-variant-base-change}{{09LU}{source}{}{sourcepointer-5}{}}
\newlabel{dga-section-variant-base-change@cref}{{[originalsection][0][]09LU}{[1][0][]source}{}{}{}}
\makeatother
\crefname{originalreference}{original reference}{original references}
\Crefname{originalreference}{Original reference}{Original references}
\makeatletter
\newlabel{lemma-ext-from-perfect-into-bounded-QCoh}{{09M4}{source}{}{sourcepointer-6}{}}
\newlabel{lemma-ext-from-perfect-into-bounded-QCoh@cref}{{[originalreference][0][]09M4}{[1][0][]source}{}{}{}}
\makeatother
\crefname{originalreference}{original reference}{original references}
\Crefname{originalreference}{Original reference}{Original references}
\makeatletter
\newlabel{lemma-tensor-with-QCoh-complex}{{09M3}{source}{}{sourcepointer-7}{}}
\newlabel{lemma-tensor-with-QCoh-complex@cref}{{[originalreference][0][]09M3}{[1][0][]source}{}{}{}}
\makeatother
\crefname{originalreference}{original reference}{original references}
\Crefname{originalreference}{Original reference}{Original references}
\makeatletter
\newlabel{proposition-compact-is-perfect}{{09M1}{source}{}{sourcepointer-8}{}}
\newlabel{proposition-compact-is-perfect@cref}{{[originalreference][0][]09M1}{[1][0][]source}{}{}{}}
\makeatother
\crefname{originalreference}{original reference}{original references}
\Crefname{originalreference}{Original reference}{Original references}
\makeatletter
\newlabel{theorem-bondal-van-den-Bergh}{{09IS}{source}{}{sourcepointer-9}{}}
\newlabel{theorem-bondal-van-den-Bergh@cref}{{[originalreference][0][]09IS}{[1][0][]source}{}{}{}}
\makeatother
\title{The derived category of quasi-coherent sheaves is a derived category of a dga}
\author{The Stacks Project Authors}
\date{stacks-project:spaces-perfect.tex:local-snapshot-20261012; GNU FDL 1.2 or later}
\begin{document}
\maketitle
\noindent\textit{Editorial note.} All delivered mathematical text is the author's own, copied from the pinned Stacks Project chapter file spaces-perfect.tex; no statement or proof step was rewritten, reordered or removed.\par\medskip
\noindent\textit{Editorial note.} The counted proof is the author's complete argument: the source gives the full assembly: generator choice, finite-Ext computation, fully faithful tensor functor with adjoint, vanishing kernel on the generator subcategory and the formal equivalence criterion\par\medskip
\noindent\textit{Editorial note.} The label theorem-DQCoh-is-Ddga names the counted statement, and the proof environment immediately following it in the source is delivered verbatim as its proof.\par\medskip
\noindent\textit{Editorial note.} The counted statement is carried by the author's own theorem environment in the chapter source.\par\medskip
\noindent\textit{Allowed standard prerequisites (editorial).}\begin{itemize}
\item the Bondal-Van den Bergh theorem producing a perfect K-injective generator of D\_QCoh(O\_X) for quasi-compact quasi-separated X
\item K-injective complexes and the identification of dga cohomology with Ext groups
\item the dga tensor functor - tensor\textasciicircum{}L\_E K and its adjoint RHom(K, -)
\item compactness of perfect objects and the fully-faithfulness criterion of the tensor functor in the compact case
\item the generator property and the formal equivalence criterion for an adjunction whose restricted kernel vanishes
\end{itemize}\medskip
\ifdefined\hypertarget\hypertarget{sourcepointer-1}{}\fi
\noindent\textit{Source reference.} Original lemma 09J1: sources/stacks-chow/tags.txt, source lines 12059--12059. Prerequisite pointer invoked at proof line(s) 4454; it names a standard fact of a sibling chapter that the counted argument uses, and every step of the counted proof built on it is retained verbatim in the delivered body.\par\medskip
\ifdefined\hypertarget\hypertarget{sourcepointer-2}{}\fi
\noindent\textit{Source reference.} Original lemma 09LZ: sources/stacks-chow/tags.txt, source lines 12162--12162. Prerequisite pointer invoked at proof line(s) 4436; it names a standard fact of a sibling chapter that the counted argument uses, and every step of the counted proof built on it is retained verbatim in the delivered body.\par\medskip
\ifdefined\hypertarget\hypertarget{sourcepointer-3}{}\fi
\noindent\textit{Source reference.} Original lemma 09LX: sources/stacks-chow/tags.txt, source lines 12160--12160. Prerequisite pointer invoked at proof line(s) 4430; it names a standard fact of a sibling chapter that the counted argument uses, and every step of the counted proof built on it is retained verbatim in the delivered body.\par\medskip
\ifdefined\hypertarget\hypertarget{sourcepointer-4}{}\fi
\noindent\textit{Source reference.} Original lemma 09LY: sources/stacks-chow/tags.txt, source lines 12161--12161. Prerequisite pointer invoked at proof line(s) 4441; it names a standard fact of a sibling chapter that the counted argument uses, and every step of the counted proof built on it is retained verbatim in the delivered body.\par\medskip
\ifdefined\hypertarget\hypertarget{sourcepointer-5}{}\fi
\noindent\textit{Source reference.} Original section 09LU: sources/stacks-chow/tags.txt, source lines 12157--12157. Prerequisite pointer invoked at proof line(s) 4414; it names a standard fact of a sibling chapter that the counted argument uses, and every step of the counted proof built on it is retained verbatim in the delivered body.\par\medskip
\ifdefined\hypertarget\hypertarget{sourcepointer-6}{}\fi
\noindent\textit{Source reference.} Original reference 09M4: sources/stacks-spaces-perfect/pinned-original-source.tex, source lines 4328--4328. Prerequisite pointer invoked at proof line(s) 4422; it names a standard fact of the same chapter that the counted argument uses, and every step of the counted proof built on it is retained verbatim in the delivered body.\par\medskip
\ifdefined\hypertarget\hypertarget{sourcepointer-7}{}\fi
\noindent\textit{Source reference.} Original reference 09M3: sources/stacks-spaces-perfect/pinned-original-source.tex, source lines 4294--4294. Prerequisite pointer invoked at proof line(s) 4444; it names a standard fact of the same chapter that the counted argument uses, and every step of the counted proof built on it is retained verbatim in the delivered body.\par\medskip
\ifdefined\hypertarget\hypertarget{sourcepointer-8}{}\fi
\noindent\textit{Source reference.} Original reference 09M1: sources/stacks-spaces-perfect/pinned-original-source.tex, source lines 4096--4096. Prerequisite pointer invoked at proof line(s) 4432; it names a standard fact of the same chapter that the counted argument uses, and every step of the counted proof built on it is retained verbatim in the delivered body.\par\medskip
\ifdefined\hypertarget\hypertarget{sourcepointer-9}{}\fi
\noindent\textit{Source reference.} Original reference 09IS: sources/stacks-spaces-perfect/pinned-original-source.tex, source lines 3891--3891. Prerequisite pointer invoked at proof line(s) 4406; it names a standard fact of the same chapter that the counted argument uses, and every step of the counted proof built on it is retained verbatim in the delivered body.\par\medskip
\begin{theorem}
\label{theorem-DQCoh-is-Ddga}
Let $S$ be a scheme.
Let $X$ be a quasi-compact and quasi-separated algebraic space over $S$.
Then there exist a differential graded algebra $(E, \text{d})$
with only a finite number of nonzero cohomology groups $H^i(E)$
such that $D_\QCoh(\mathcal{O}_X)$ is equivalent
to $D(E, \text{d})$.
\end{theorem}

\begin{proof}
Let $K^\bullet$ be a K-injective complex of $\mathcal{O}$-modules which
is perfect and generates $D_\QCoh(\mathcal{O}_X)$. Such a
thing exists by Theorem \ref{theorem-bondal-van-den-Bergh}
and the existence of K-injective resolutions. We will show the
theorem holds with
$$
(E, \text{d}) = \Hom_{\text{Comp}^{dg}(\mathcal{O}_X)}(K^\bullet, K^\bullet)
$$
where $\text{Comp}^{dg}(\mathcal{O}_X)$ is the differential graded category
of complexes of $\mathcal{O}$-modules. Please see
Differential Graded Algebra, Section \ref{dga-section-variant-base-change}.
Since $K^\bullet$ is K-injective we
have
\begin{equation}
\label{equation-E-is-OK}
H^n(E) = \Ext^n_{D(\mathcal{O}_X)}(K^\bullet, K^\bullet)
\end{equation}
for all $n \in \mathbf{Z}$. Only a finite number of these Exts
are nonzero by Lemma \ref{lemma-ext-from-perfect-into-bounded-QCoh}.
Consider the functor
$$
- \otimes_E^\mathbf{L} K^\bullet :
D(E, \text{d}) \longrightarrow D(\mathcal{O}_X)
$$
of
Differential Graded Algebra, Lemma
\ref{dga-lemma-tensor-with-complex-derived}.
Since $K^\bullet$ is perfect, it defines a compact object of
$D(\mathcal{O}_X)$, see Proposition \ref{proposition-compact-is-perfect}.
Combined with (\ref{equation-E-is-OK}) the functor above is fully
faithful as follows from
Differential Graded Algebra, Lemmas
\ref{dga-lemma-fully-faithful-in-compact-case}. It has a right adjoint
$$
R\Hom(K^\bullet, - ) : D(\mathcal{O}_X) \longrightarrow D(E, \text{d})
$$
by Differential Graded Algebra, Lemmas
\ref{dga-lemma-tensor-with-complex-hom-adjoint}
which is a left quasi-inverse functor by generalities on adjoint
functors. On the other hand, it follows from
Lemma \ref{lemma-tensor-with-QCoh-complex} that we obtain
$$
- \otimes_E^\mathbf{L} K^\bullet :
D(E, \text{d}) \longrightarrow D_\QCoh(\mathcal{O}_X)
$$
and by our choice of $K^\bullet$ as a generator of
$D_\QCoh(\mathcal{O}_X)$ the kernel of the adjoint
restricted to $D_\QCoh(\mathcal{O}_X)$ is zero.
A formal argument shows that we obtain the desired equivalence, see
Derived Categories, Lemma
\ref{derived-lemma-fully-faithful-adjoint-kernel-zero}.
\end{proof}

\end{document}
